% id: 205-08
% book: 베이직쎈 대수 · 12 수학적 귀납법
% section: 실전 감각 UP
% figure: none
% answer: 
% answer_source: 
% variant_level: 0 (원본 전사)
% 이 파일은 build.mjs 가 items.json 에서 만든다. 고칠 때는 items.json 을 고친다.
\smash{\raisebox{4.5mm}{\dmkichul{베이직쎈 대수 205쪽 08번 · 꼭 나와요}}}
다음은 $n\ge 2$인 모든 자연수 $n$에 대하여 \par\smallskip\dispeq{$1+\dfrac{1}{2^2}+\dfrac{1}{3^2}+\cdots+\dfrac{1}{n^2}<2-\dfrac{1}{n}$}\par\smallskip\noindent 이 성립함을 수학적 귀납법으로 증명한 것이다.\exprbox{\begin{minipage}{0.96\linewidth}\raggedright\lineskip=4pt {\small\bfseries 증명}\par\smallskip (i) $n=2$일 때,\\ \hspace*{1.5em}(좌변)$=1+\dfrac{1}{2^2}=\dfrac{5}{4}$, (우변)$=2-\dfrac{1}{2}=\dfrac{3}{2}$\\ \hspace*{1em}따라서 주어진 부등식이 성립한다.\\ (ii) $n=k\ (k\ge 2)$일 때 주어진 부등식이 성립한다고 가정하면\\ \hspace*{1.5em}$1+\dfrac{1}{2^2}+\dfrac{1}{3^2}+\cdots+\dfrac{1}{k^2}<2-\dfrac{1}{k}$\\ \hspace*{1em}위의 식의 양변에 $\dfrac{1}{(k+1)^2}$을 더하면\\ \hspace*{1.5em}$1+\dfrac{1}{2^2}+\dfrac{1}{3^2}+\cdots+\dfrac{1}{k^2}+\dfrac{1}{(k+1)^2}$\\ \hspace*{1.5em}$<2-\dfrac{1}{k}+\dfrac{1}{(k+1)^2}$\\ \hspace*{1em}이때\\ \hspace*{1.5em}$\left\{2-\dfrac{1}{k}+\dfrac{1}{(k+1)^2}\right\}-\left(2-\framebox[2.2em]{\scriptsize ㈎}\right)$\\ \hspace*{1.5em}$=-\dfrac{1}{k}+\dfrac{1}{(k+1)^2}+\dfrac{1}{k+1}$\\ \hspace*{1.5em}$=-\dfrac{1}{k(k+1)^2}<\framebox[2.2em]{\scriptsize ㈏}$\\ \hspace*{1em}이므로 $2-\dfrac{1}{k}+\dfrac{1}{(k+1)^2}<2-\framebox[2.2em]{\scriptsize ㈎}$\\ \hspace*{1.5em}$\therefore 1+\dfrac{1}{2^2}+\dfrac{1}{3^2}+\cdots+\dfrac{1}{k^2}+\dfrac{1}{(k+1)^2}$\\ \hspace*{1.5em}$<2-\framebox[2.2em]{\scriptsize ㈎}$\\ \hspace*{1em}따라서 $n=k+1$일 때도 주어진 부등식이 성립한다.\\ (i), (ii)에서 $n\ge 2$인 모든 자연수 $n$에 대하여 부등식이 성립한다.\end{minipage}}\noindent 위의 과정에서 ㈎, ㈏에 알맞은 것을 각각 $f(k)$, $a$라 할 때, $f(a)$의 값은?
\dmchoicesauto{}{$-1$}{$0$}{$1$}{$2$}{$3$}
